% ============================================================
%  Section 2 --- Formalisation
% ============================================================

\section{Formalisation du problème}
\label{sec:formalisation}

Une formalisation explicite est nécessaire pour deux raisons. D'abord, elle contraint les
choix d'implémentation : les règles d'agrégation intra-service, le nombre de régimes $|\Theta|$
et le budget de latence découlent directement des définitions formelles.
Ensuite, elle définit les critères de falsification des hypothèses de manière non ambiguë,
ce qui rend les résultats négatifs aussi informatifs que les résultats positifs.

\subsection{Graphe de services}
\label{sec:graph}

Soit $G(t) = (V, E(t), w_E(t))$ le graphe de services à l'instant $t$, où $V$ est
l'ensemble des services Kubernetes (Deployments), constant avec $|V| = N = 6$,
et $E(t)$ les arêtes pondérées par
$w_E : E(t) \to \mathbb{R}^3$, $e_{ij}(t) \mapsto (\text{volume}_{ij},\,
\text{latence\_med}_{ij},\, \text{taux\_erreur}_{ij})$.

Une arête $e_{ij}$ est présente à l'instant $t$ si et seulement si le volume de requêtes
entre le service $i$ et le service $j$ est strictement positif sur la fenêtre glissante
précédente. Les poids sont agrégés depuis les spans OTel selon les règles suivantes,
choisies pour respecter les propriétés statistiques de chaque type de mesure :
\begin{itemize}
  \item Saturation (CPU, RAM, \texttt{net\_sat}, \texttt{disk\_io}) : \textbf{maximum}
    sur les pods du service --- la saturation la plus haute détermine le comportement.
  \item Taux (\texttt{error\_rate}, \texttt{log\_warn\_rate}) : \textbf{somme pondérée
    par volume} --- un service à faible trafic et taux élevé ne doit pas être dilué.
  \item Latence (\texttt{latency\_p99}, \texttt{span\_dur\_median}) : \textbf{percentile
    99 sur l'union} des distributions individuelles des pods --- jamais un percentile de
    percentiles, qui sous-estime les queues de distribution \cite{gregg2013}.
  \item Métriques structurelles (\texttt{trace\_depth}, \texttt{fan\_out},
    \texttt{lexical\_entropy}) : \textbf{médiane} --- robuste aux outliers transitoires.
\end{itemize}

\subsection{Signal de télémétrie}
\label{sec:signal}

Le signal $S(t) \in \mathbb{R}^{N \times 17}$ est défini comme la concaténation de
trois modalités :
\[
  S(t) = \bigl[M(t) \mid T(t) \mid L(t)\bigr]
\]
avec :
\begin{itemize}
  \item $M(t) \in \mathbb{R}^{N \times 7}$ --- métriques Prometheus, couvrant les ressources
    selon la méthodologie USE \cite{gregg2013} (utilisation, saturation, erreurs) :
    \texttt{cpu\_util}, \texttt{ram\_util}, \texttt{latency\_p99},
    \texttt{error\_rate\_http}, \texttt{net\_sat}, \texttt{disk\_io},
    \texttt{queue\_depth}.
  \item $T(t) \in \mathbb{R}^{N \times 6}$ --- traces OTel, capturant la structure et
    la qualité des appels inter-services :
    \texttt{span\_dur\_median}, \texttt{abnormal\_span\_rate}, \texttt{trace\_depth},
    \texttt{fan\_out}, \texttt{retry\_rate}, \texttt{latency\_cv}.
  \item $L(t) \in \mathbb{R}^{N \times 4}$ --- logs OTel, dont deux features sémantiques :
    \texttt{log\_error\_rate}, \texttt{log\_warn\_rate},
    \texttt{semantic\_anomaly} (distance SentenceBERT au centroïde normal),
    \texttt{lexical\_entropy}.
\end{itemize}

La modalité $M(t)$ fournit les signaux de ressource matures et stables ; $T(t)$ capture
les effets de propagation dans le graphe d'appel ; $L(t)$ enrichit avec une sémantique
applicative non disponible dans les métriques brutes.
L'ablation par réentraînement (Section~\ref{sec:ablation_h1}) montre que $M(t)$ seul suffit
pour la structuration géométrique (H1), tandis que $T(t)$ et $L(t)$ sont indispensables
pour la prédictibilité des précurseurs (H3).

\subsection{Régimes opérationnels}
\label{sec:regimes}

Nous définissons quatre régimes $\theta(t) \in \Theta$ :
\[
  \Theta = \{\theta_\text{normal},\; \theta_\text{drift},\;
             \theta_\text{anomaly},\; \theta_{\text{drift} \cap \text{anomaly}}\}
\]

$\theta_\text{normal}$ est l'état nominal du système, sans injection.
$\theta_\text{drift}$ correspond aux changements bénins de distribution --- déploiements
progressifs, autoscaling, modifications de quota --- qui ne dégradent pas les SLOs.
$\theta_\text{anomaly}$ regroupe les vrais incidents : saturation ressource, crash, fuite
mémoire, latence anormale.
$\theta_{\text{drift} \cap \text{anomaly}}$ modélise le cas hybride d'un déploiement
défectueux : la distribution change (drift) et les SLOs se dégradent simultanément (anomalie).
Ce quatrième régime exige un mécanisme de look-through dans l'étape de détection de drift,
afin de ne pas supprimer le signal d'anomalie sous prétexte qu'un drift est en cours.
Trois régimes seraient insuffisants : le régime $\theta_{\text{drift} \cap \text{anomaly}}$
ne peut pas être correctement géré ni comme drift pur (l'alerte serait supprimée)
ni comme anomalie pure (la logique de recalibration du détecteur serait incorrecte).

\subsection{Pipeline EWAT}
\label{sec:pipeline_overview}

Le pipeline se décompose en quatre étapes séquentielles, conçues pour être indépendantes
et testables unitairement :

\begin{description}
  \item[\textbf{Étape 0} --- Détection de drift (MMD-RFF).]
    $\widehat{\text{MMD}}^2(W_\text{ref}, W_\text{cur})$ via Random Fourier Features
    avec mécanisme \emph{look-through} \cite{hinder2024}.
    La complexité est $O(nD)$ avec $D = 256$ features aléatoires, permettant une inférence
    en temps réel. Le seuil $\varepsilon_\text{drift} = 0.5226$ est calibré par critère
    de Youden sur les drifts bénins injectés (AUC $= 0.60$ sur train).
    Budget de latence : $< 1$\,s.

  \item[\textbf{Étape 1} --- Encodeur STGCN.]
    $z_e = \text{Enc}_\theta(S_{[t-W,t+\delta]}, G(t)) \in \mathbb{R}^{64}$,
    convolution spatio-temporelle sur le graphe de services \cite{yu2018}.
    L'encodeur joint la dimension spatiale (structure du graphe $G(t)$) et la dimension
    temporelle (série $S(t)$) via un réseau TCN, produisant un embedding
    $z_e \in \mathbb{R}^{64}$ par épisode.
    Budget : $< 2$\,s.

  \item[\textbf{Étape 2} --- Typage contrastif.]
    Réseau siamois appris sur les paires de scénarios Chaos Mesh :
    $d_\varphi(z_i, z_j) \to 0$ si même type, $\to 1$ sinon.
    Le clustering hiérarchique agglomératif sur les embeddings produit
    $K$ types $\{C_1, \ldots, C_K\}$ sélectionnés par score de silhouette sur train
    ($K = 10$ sur ewat\_v3).
    En validation et test, les labels sont assignés par plus proche centroïde depuis les
    centroides train --- ce qui garantit l'absence de fuite d'information inter-splits.

  \item[\textbf{Étape 2b} --- Ontologie.]
    Relations temporelles ($C_i \to C_j$, support $\geq 3$),
    causales (Transfer Entropy, estimateur KSG \cite{kraskov2004},
    100 permutations, BH-FDR \cite{benjamini1995})
    et de co-occurrence ($\chi^2$ Yates, correction Holm--Bonferroni).

  \item[\textbf{Étape 3} --- Précurseurs typés.]
    $\hat{p}_i(t) = f_i(S_{[t-k,t]}, G(t)) \in [0,1]$,
    $k^*_i = \arg\max_k \text{AUROC}(f_i, k)$ sélectionné sur val,
    évalué sur test uniquement.
    Budget : $< 1$\,s.
\end{description}

La sortie est $\text{Alert}(t) = (C_i,\, \hat{p}_i(t),\, k^*_i,\, \text{fiche}_{C_i})$,
enrichie d'une fiche d'interprétabilité par cluster générée par permutation importance
et validée par KernelSHAP.

\subsection{Hypothèses et critères de falsification}
\label{sec:hypotheses}

Chaque hypothèse est formulée avec un critère quantitatif de falsification, de sorte
qu'un résultat négatif soit aussi informatif qu'un résultat positif.

\textbf{H1 --- Structurabilité.}
Les embeddings STGCN forment des clusters exploitables dans l'espace latent,
définis par silhouette $\geq 0.3$ en held-out sur 5 graines $\times$ 5 splits.
Le seuil $0.3$ est justifié par Kaufman \& Rousseeuw \cite{kaufman1990} comme limite
minimale d'une structure exploitable.

\textbf{H2a --- Séparabilité drift par look-through.}
Le mécanisme look-through réduit significativement le FPR sur anomalie à rappel constant,
mesuré par un test de Student unilatéral apparié ($p < 0.05$) sur les 45 épisodes test.

\textbf{H2b --- Identification du régime $\theta_{\text{drift} \cap \text{anomaly}}$.}
Le cluster C8 (\texttt{faulty\_deploy\_overlap}) présente un overlap drift\,+\,alerte
significativement supérieur aux clusters de drift pur (C5, C6, C9),
testé par Fisher exact ou proportion bootstrap.

\textbf{H3 --- Prédictibilité des précurseurs typés.}
L'AUROC par type est supérieur à la baseline aléatoire ($0.5$) pour chaque horizon $k$,
$k^*_i$ sélectionné sur val et évalué sur test.
Les intervalles de confiance à 95\,\% sont estimés par bootstrap (1000 rééchantillonnages).
